\section[Discrete physical realization on TPK routes]{Discrete physical realization:
positive and oriented readings on TPK routes, common domain, and
reversible localization}
\label{sec:int68-realizacion-fisica-discreta}

The spectral form, the Hilbert realization, the dimensional section, and Cartan transport are not decorations appended to a classical output. They are four coordinated representations of the same enriched TPK route. To express that causal unity, their common domain must first be constructed; it is not sufficient to write four arrows with the same source symbol.

\subsection{Common domain as a fiber product}

Let \(\mathsf P_{\rm TPK}^{\rm enr}\) be the category of prepared TPK states
and registered routes compatible with their endpoints, boundary, and restrictions.
Let
\[
 \mathsf D_{\rm sp}:=\mathsf{SurvPath}_{\rm HMT},\qquad
 \mathsf D_Q:=\operatorname{Dom}(\mathcal Q),\qquad
 \mathsf D_D:=\operatorname{Dom}(\mathfrak D_{\rm HMT}),\qquad
 \mathsf D_C:=\operatorname{Dom}(\rho_C^{\rm disc}),
\]
with their faithful forgetful functors
\[
 U_i:\mathsf D_i\longrightarrow\mathsf P_{\rm TPK}^{\rm enr}
 \qquad(i\in\{{\rm sp},Q,D,C\}).
\]
Each \(U_i\) forgets only the realization and preserves the underlying TPK state and route.
We define
\begin{equation}
 \mathsf P_{\rm CT}^{+}
 :=
 \mathsf D_{\rm sp}
 \times_{\mathsf P_{\rm TPK}^{\rm enr}}\mathsf D_Q
 \times_{\mathsf P_{\rm TPK}^{\rm enr}}\mathsf D_D
 \times_{\mathsf P_{\rm TPK}^{\rm enr}}\mathsf D_C.
 \label{eq:int68-pct-pullback}
\end{equation}
An object is therefore a prepared state equipped with the four
realizations over the same antecedent; an arrow is a recorded route
whose underlying image coincides in the four factors. If the
presentations are identified only up to isomorphism, the corresponding bicategorical
fiber product is used; in the fixed presentation of the record
the strict fiber product of
\eqref{eq:int68-pct-pullback} is adopted.

\begin{proposition}[Existence of an object in the common domain]
\label{prop:int68-pct-no-vacio}
The category \(\mathsf P_{\rm CT}^{+}\) is nonempty. It is not necessary
to specialize the central electronic route in order to prove this: any certified
coinductive section, equipped with its identity arrow and its neutral data, admits
compatible lifts to the four domains of
\eqref{eq:int68-pct-pullback}.
\end{proposition}

\begin{proof}
The \cref{thm:generacion-monodromica-conjunta-rev7} provides, for example,
a prepared surviving section \(x_\pi\). Its identity
\(\operatorname{id}_{x_\pi}\) is simultaneously: a surviving route;
a morphism of \(\mathsf P_{\rm TPK}\), which \(\mathcal Q\) sends to the
bounded identity of \(\mathcal Q(x_\pi)\); the empty route with null signature,
which the \cref{thm:funtor-dimensional-minimo} sends to
\((0,0,0,1)\); and the identity of the groupoid, which the
\cref{thm:realizacion-discreta-cartan-rev6} sends to the neutral Poincare element
\((I,0)\). These four lifts forget to the same object
\(x_\pi\) and the same morphism \(\operatorname{id}_{x_\pi}\). Their quadruple
defines, by the fiber-product property, an object of
\(\mathsf P_{\rm CT}^{+}\).
\end{proof}

\begin{definition}[Adapters to the common domain]
\label{def:int68-pct-positive}
The four adapters are the canonical projections
\[
 \begin{aligned}
 F_{\rm sp}:\mathsf P_{\rm CT}^{+}&\longrightarrow
     \mathsf{SurvPath}_{\rm HMT},\\
 F_Q:\mathsf P_{\rm CT}^{+}&\longrightarrow\operatorname{Dom}(\mathcal Q),\\
 F_D:\mathsf P_{\rm CT}^{+}&\longrightarrow
     \operatorname{Dom}(\mathfrak D_{\rm HMT}),\\
 F_C:\mathsf P_{\rm CT}^{+}&\longrightarrow
     \operatorname{Dom}(\rho_C^{\rm disc}).
 \end{aligned}
\]
In particular, \(F_{\rm sp}\) preserves the compatible family of
finite restrictions of the section and the route; it does not reconstruct that family
from its spectrum.
\end{definition}

\begin{lemma}[Functoriality of the adapters]
\label{lem:int68-adaptadores}
The four adapters preserve identities and composition.
\end{lemma}

\begin{proof}
The identity of an object of the fibered product is the quadruple of identities
on the same TPK state; each projection returns the identity of the factor.
Composition is defined componentwise between quadruples that
project onto the same concatenation of routes. Projecting before or after
composition gives, by definition, the same result.
\end{proof}

\subsection{\(4\) Components and Positive Product}

The spectral, Hilbertian, dimensional, and Cartan realization theorems provide
\[
 \mathfrak S:\mathsf{SurvPath}_{\rm HMT}
      \longrightarrow\mathbf{SpecRoute},
 \qquad
 \mathcal Q:\operatorname{Dom}(\mathcal Q)
      \longrightarrow\mathsf{Hilb}_{\mathbb R,J},
\]
\[
 \mathfrak D_{\rm HMT}:\operatorname{Dom}(\mathfrak D_{\rm HMT})
      \longrightarrow\mathbf B\mathfrak M_D^+,
 \qquad
 \rho_C^{\rm disc}:\operatorname{Dom}(\rho_C^{\rm disc})
      \longrightarrow\mathbf B G_C,
\]
where
\[
 G_C:=\operatorname{Spin}^{+}(1,3)\ltimes\mathbb R^{1,3}
\]
and
\begin{equation}
 \mathfrak M_D^+
 :=\mathbb N\times\mathbb R_{\geq0}^{2}\times\mathbb R_{>0},
 \qquad
 (n,a,b,\lambda)\odot(n',a',b',\lambda')
 =(n+n',a+a',b+b',\lambda\lambda').
 \label{eq:int68-monoide-dimensional-positivo}
\end{equation}
Here \(\mathbf BM\) is the one-object category associated with the monoid
\(M\). The first functor is that of
\cref{thm:functor-realizacion-espectral-numero}; the second is the
Hilbert realization defined in the postulates chapter; the third is
\cref{thm:funtor-dimensional-minimo}; and the fourth is
\cref{thm:realizacion-discreta-cartan-rev6}.

We define the common constraints
\[
 \mathfrak S_+:=\mathfrak S\circ F_{\rm sp},\qquad
 \mathcal Q_+:=\mathcal Q\circ F_Q,\qquad
 \mathfrak D_+:=\mathfrak D_{\rm HMT}\circ F_D,\qquad
 \rho_{C,+}^{\rm disc}:=\rho_C^{\rm disc}\circ F_C.
\]

\begin{theorem}[Positive discrete realization on the common domain]
\label{thm:int68-producto-positivo}
The map
\begin{equation}
 \mathfrak R_{\rm MD}^{\rm disc,+}
 :=(\mathfrak S_+,\mathcal Q_+,\mathfrak D_+,
       \rho_{C,+}^{\rm disc})
 \label{eq:int68-producto-positivo}
\end{equation}
is a functor
\[
 \mathsf P_{\rm CT}^{+}\longrightarrow
 \mathbf{SpecRoute}\times\mathsf{Hilb}_{\mathbb R,J}
 \times\mathbf B\mathfrak M_D^+\times\mathbf BG_C.
\]
It is not asserted that every TPK route belongs to \(\mathsf P_{\rm CT}^{+}\).
\end{theorem}

\begin{proof}
By \cref{lem:int68-adaptadores}, each projection \(F_i\) is a functor;
by the cited theorems, so are
\(\mathfrak S,\mathcal Q,\mathfrak D_{\rm HMT}\) and
\(\rho_C^{\rm disc}\) on their declared domains. Their compositions
therefore preserve identities and concatenation. Composition in the product
codomain is computed componentwise, which proves
\eqref{eq:int68-producto-positivo}.
\end{proof}

The Hilbert component has general target
\(\mathsf{Hilb}_{\mathbb R,J}\), whose morphisms are bounded real
operators intertwining the complex structures. Let
\(\mathsf P_{\rm CT}^{u}\subseteq\mathsf P_{\rm CT}^{+}\) be the subcategory
of arrows \(\gamma:x\to y\) for which
\begin{equation}
 \mathcal Q_+(\gamma)^*\mathcal Q_+(\gamma)=I_{\mathcal Q_+(x)},
 \qquad
 \mathcal Q_+(\gamma)\mathcal Q_+(\gamma)^*=I_{\mathcal Q_+(y)}.
 \label{eq:int68-unitariedad-doble}
\end{equation}
The identities and products of unitary operators are again
unitary; hence this condition defines a subcategory, and only within it is
obtained
\begin{equation}
 \mathcal U_{\rm reg}:=
 \mathcal Q_+\big|_{\mathsf P_{\rm CT}^{u}}:
 \mathsf P_{\rm CT}^{u}\longrightarrow
 \mathsf{UHilb}_{\mathbb R,J}.
 \label{eq:int68-restriccion-unitaria}
\end{equation}
One-sided invertibility or norm conservation on a subset does not
replace the two identities of \eqref{eq:int68-unitariedad-doble}.

The seven-entry notation decomposes into the four typed components: \(Q\) is \(\mathcal Q_+\); \(U\) is the restriction \eqref{eq:int68-restriccion-unitaria}; \(A\) is the action coordinate of \(\mathfrak D_+\); \(C\) is the conjugation declared on the subdomain where it intertwines the Hilbert component and orientation; \(\mathcal E\) and \(\mathcal B\) are constitutive readers subordinated to the dimensional data; and \(N_{\rm sp}\) is a quadratic property of the split realization. None of the final four is added to the product as a fifth functor without a domain and codomain.

\subsection{Positive and oriented readings}

Let \(e\mapsto e^\vee\) be edge reversal in the underlying directed
graph. In \(\mathsf P_{\rm CT}^{+}\), \(e^\vee\) is another edge and is not
yet declared the categorical inverse of \(e\). Having fixed a weight
\(w(e)=w(e^\vee)>0\), we define on a route
\(\gamma=e_1\cdots e_r\)
\[
 V_+(\gamma):=\sum_{j=1}^r w(e_j),
 \qquad
 V_{\rm or}(\gamma):=
 \sum_{j=1}^r\varepsilon(e_j)w(\underline e_j),
\]
where \(\underline e_j\) is the unoriented edge and
\(\varepsilon(e^\vee)=-\varepsilon(e)\). The first reader preserves
total positive variation; the second, net oriented transport. Both are
additive under concatenation and receive the same TPK route.

\begin{proposition}[No factorization of the positive reader by the oriented]
\label{prop:int68-no-factorizacion}
Suppose that the identity and the loop
\(ee^\vee:=e^\vee\circ e\) belong to the constructed domain. There is no
function \(H\) on \(\operatorname{im}V_{\rm or}\) such that
\(V_+=H\circ V_{\rm or}\).
\end{proposition}

\begin{proof}
We have
\[
 V_{\rm or}(\operatorname{id})=0
 =V_{\rm or}(ee^\vee),
 \qquad
 V_+(\operatorname{id})=0,
 \qquad
 V_+(ee^\vee)=2w(e)>0.
\]
The same input \(0\) would force \(H(0)\) to take two distinct values.
\end{proof}

This result is a concrete nonfactorization in the constructed domain; it
is not a universal theorem about every possible enriched intertwiner.
In particular, it does not authorize erasing a reader. The unit lies in the
common TPK antecedent and not in a horizontal compression between its outputs.

\subsection{Reversible localization and oriented realization}

Let \(\mathsf P_{{\rm CT},{\rm rev}}^+\) be the subcategory generated by the
edges for which the following four conditions have been verified:
\begin{enumerate}
 \item \(\mathfrak S_+(e)\) is an isomorphism of
       \(\mathbf{SpecRoute}\) and
       \(\mathfrak S_+(e^\vee)=\mathfrak S_+(e)^{-1}\);
 \item \(\mathcal Q_+(e)\) satisfies
       \eqref{eq:int68-unitariedad-doble} and
       \(\mathcal Q_+(e^\vee)=\mathcal Q_+(e)^*\);
 \item the oriented dimensional datum arises from an additive cochain with
       \(d_{\rm or}(e^\vee)=-d_{\rm or}(e)\) and takes values in
       \(\operatorname{Gr}(\mathfrak M_D^+)\);
 \item the Cartan cocycles are normalized, additive, and antisymmetric
       under reversal.
\end{enumerate}
We form the groupoid
\[
 \Pi_\vee(\mathsf P_{{\rm CT},{\rm rev}}^+)
\]
obtained from the groupoid completion by additionally imposing
\([e^\vee]=[e]^{-1}\). The oriented dimensional component takes values in
\begin{equation}
 \operatorname{Gr}(\mathfrak M_D^+)
 \cong\mathbb Z\times\mathbb R^2\times\mathbb R_{>0},
 \label{eq:int68-grothendieck-dimensional}
\end{equation}
with inverse
\((n,a,b,\lambda)^{-1}=(-n,-a,-b,\lambda^{-1})\). This signed reader is not
the positive magnitude \(V_+\) relabeled: its values on the reverse
edges are distinct.

For the Cartan component, let \(c_g,c_s\) be cocycles as in item 4,
let \(R_4\in\operatorname{Spin}^{+}(1,3)\), let \(R_4u=u\), and let
\(\ell_0>0\). Then
\[
 \rho_C^{\rm disc}(\gamma)
 =\bigl(R_4^{c_g(\gamma)},\ell_0c_s(\gamma)u\bigr).
\]
The law
\((R,a)(S,b)=(RS,a+Rb)\), the equality \(R_4u=u\), and the additivity of the
cocycles give
\[
 \rho_C^{\rm disc}(\gamma_2\gamma_1)
 =\rho_C^{\rm disc}(\gamma_2)\rho_C^{\rm disc}(\gamma_1),
 \qquad
 \rho_C^{\rm disc}(\gamma^\vee)
 =\rho_C^{\rm disc}(\gamma)^{-1}.
\]

\begin{theorem}[Discrete oriented realization]
\label{thm:int68-realizacion-orientada}
Under the four preceding conditions, the components descend uniquely
to a functor
\[
 \mathfrak R_{\rm MD}^{\rm disc,or}:
 \Pi_\vee(\mathsf P_{{\rm CT},{\rm rev}}^+)
 \longrightarrow
 \operatorname{Core}(\mathbf{SpecRoute})
 \times\mathsf{UHilb}_{\mathbb R,J}
 \times\mathbf B\operatorname{Gr}(\mathfrak M_D^+)
 \times\mathbf BG_C.
\]
In particular,
\[
 \mathfrak R_{\rm MD}^{\rm disc,or}(\gamma^{-1})
 =\mathfrak R_{\rm MD}^{\rm disc,or}(\gamma)^{-1}.
\]
\end{theorem}

\begin{proof}
Conditions 1 and 2 make the spectral and
Hilbert images invertible; condition 3 takes values in a group; and the Cartan
component already takes values in \(G_C\). All four respect the generating relation
\([e^\vee]=[e]^{-1}\). By the universal property of groupoid
localization, each descends uniquely; their product is the indicated
functor. The inversion identity is verified componentwise.
\end{proof}

The positive reader \(\mathfrak R_{\rm MD}^{\rm disc,+}\) and the oriented reader
\(\mathfrak R_{\rm MD}^{\rm disc,or}\) are siblings: neither is defined as a
compression of the other. The subsequent passage from an oriented value to its modulus is
not, in general, monoidal and coincides with a positive reading only on a cone
without cancellation. Likewise, this pair does not replace the foundational double projection:
the Archimedean evaluation and the
Paley--Witt--Golay--Leech--\(V^\natural\)--Monster--Moonshine branch remain
sibling readers of the same enriched state. The four components
constructed in this section replace neither of those two projections.

\subsection{Cartan Entanglement and Soft Limit}

Let
\[
 \mathcal A_{\ell_0}
 =\begin{pmatrix}\omega&\ell_0^{-1}e\\0&0\end{pmatrix}
\]
a Cartan connection declared. The equation
\begin{equation}
 \operatorname{Hol}_{\mathcal A_{\ell_0}}
   \bigl(\mathfrak R_C(\gamma)\bigr)
 =\rho_C\bigl(\operatorname{Hol}_{\rm TPK}(\gamma)\bigr)
 \label{eq:int68-cuadrado-cartan}
\end{equation}
is exact in the discrete model when both holonomies are defined by the same
ordered product of edges.

\begin{criterion}[Promotion to a smooth realization]
\label{crit:int68-cartan-suave}
The equation \eqref{eq:int68-cuadrado-cartan} defines a smooth prolongation only
if subdivision independence, convergence of the ordered
product, and local recovery of a co-tetrad and a connection with
the declared domains and regularity are also proved.
\end{criterion}

The term ‘physical’ means here that the four components publish
spectrum, state realization, dimensions, and Cartan transport without
losing their TPK antecedent\@. It does not assert uniqueness, equivalence of categories,
coverage of every route, an automatic SI chart, or a complete field theory.

\begin{criterion}[Obstructions to positive, unitary, and oriented realizations]
\label{crit:int68-obstrucciones-realizacion}
The positive realization is not defined if any of the four
adapters is absent, if their images do not have the same underlying TPK route, or if a
component does not preserve identity and composition. The unitary restriction does not
exist if either of the two equalities in
\eqref{eq:int68-unitariedad-doble} fails. The oriented realization does not descend to the
groupoid if a spectral image is not isomorphic, if a localized relation
changes its image, if a positive magnitude is treated as an inverse, if a
reversal does not change the sign of the oriented data, or if the Cartan
cocycle ceases to be normalized, additive, or antisymmetric. No compatible
smooth prolongation exists if the result depends on the subdivision or if
the limit does not recover the cotetrad and connection. Any use of an external
coordinate to select the route violates the constructed domain.
\end{criterion}

\paragraph{Domain, codomain, and status of the HMT–MD confluent equations}
The positive product is proved on the common fiber product, and the
oriented realization on reversible localization under the four
explicit conditions. Extending the domain to other routes requires constructing their
four lifts over a common antecedent. Smooth prolongation additionally requires
the subdivision, convergence, and local recovery properties of
\cref{crit:int68-cartan-suave}; a finite verification does not replace those
properties.
